% !TeX spellcheck = de_DE

\subsection{Grundlagen zu Large Language Models (LLMs)}
Large Language Models (LLMs) stellen einen spezialisierten Typ von Modellen der künstlichen Intelligenz dar, der darauf trainiert wurde, menschenähnliche Texte zu verstehen und zu generieren.
LLMs basieren auf der statistischen Erfassung sprachlicher Muster durch die Verarbeitung großer Textdatensätze und zeichnen sich in der Regel durch eine enorme Kapazität von Milliarden von Parametern aus.
Sie bilden den Kern einer neuen Ära von sogenannten Foundation Models, Basismodelle, die auf breiten Datenmengen trainiert werden und die Fähigkeit besitzen, an eine Vielzahl unterschiedlicher Aufgaben angepasst zu werden \cite[S. 1, 3]{bommasani_opportunities_2022}.\\
Die technologische Grundlage von fast allen modernen LLMs ist die Transformer-Architektur.
Im Gegensatz zu älteren neuronalen Netzen ermöglichen Transformer eine parallele Verarbeitung von Eingaben, was ihre Skalierbarkeit stark erhöht.
Das Herzstück dieser Architektur bildet der Self-Attention-Mechanismus.
Dieser erlaubt es dem Modell, jedes Wort innerhalb einer Sequenz im Verhältnis zu allen anderen Wörtern zu gewichten und somit kontextuelle Zusammenhänge und Beziehungen unabhängig von der Distanz der Wörter zu erfassen \cite[S. 1-2]{vaswani_attention_nodate}.
Durch den Einsatz von Multi-Head-Attention kann das Modell zudem gleichzeitig verschiedene Arten von Beziehungen und Mustern an unterschiedlichen Positionen der Eingabe erkennen.
Da Transformer keine Informationen über die Reihenfolge von Token besitzen, werden diese durch Positional Encoding (Positionskodierung) explizit in die Eingabe-Embeddings injiziert \cite[S. 4-6]{vaswani_attention_nodate}.

Der Entwicklungsprozess dieser Modelle folgt hauptsächlich dem Ablauf von Pre-training und Fine-tuning.
Im Rahmen des Pre-trainings lernen die LLMs durch selbst überwachtes Lernen, indem sie beispielsweise fehlende Wörter in Datensätzen von zum Beispiel Wikipedia-Dateien vorhersagen.
Bevor die Verarbeitung erfolgt, wird der Text mittels Tokenisierung in kleine Einheiten unterteilt, um auch Wörter außerhalb des festen Vokabulars zu haben.
Nach dem allgemeinen Vortraining kann das Modell durch Fine-tuning mit kleineren, aufgabenspezifischen Datensätzen für konkrete Anwendungen wie Übersetzung oder auch Leichte Sprache spezialisiert werden.

Ein wesentliches Merkmal von LLMs ist die Emergenz, womit Fähigkeiten gemeint sind, die erst durch eine massive Erhöhung der Modellgröße, Datenmenge und Rechenleistung entstehen.
Ein Beispiel für Emergenz ist das In-Context-Learning.
Hierbei kann das Modell neue Aufgaben allein durch die Bereitstellung von Anweisungen oder Beispielen innerhalb des Eingabetextes(Prompts) lösen, ohne dass das Modell auf die Aufgabe spezifisch angepasst werden muss.
In der Forschung wird hierbei zwischen Zero-shot(Prompts ohne Beispiele), One-shot(ein Beispiel) und Few-shot(einige wenige Beispiele) unterschieden.
Die Leistung dieser Modelle folgt dabei oft einem Potenzgesetz in Abhängigkeit von Rechenaufwand und Modellkapazität.

Trotz der beeindruckenden Fähigkeiten von LLMs, haben sie auch einige Probleme.
Sie sind anfällig für Halluzinationen, wobei falsche Fakten oder unsinnige Antworten mit hoher Zuversicht von den Modellen generiert werden.
Zudem können sie soziale Vorurteile widerspiegeln, die in den nicht überwachten Internet-Trainingsdaten enthalten sind.
Ein weiteres Problem ist die Datenkontamination, dabei werden Testdaten von Benchmarks unbeabsichtigt Teil des Trainings und können so die Ergebnisse verfälschen.
Um Halluzinationen zu reduzieren, werden LLMs zunehmend in Retrival-Argumented Generation(RAG)-Systeme integriert, dabei wird das generative Modell mit externen, verifizierbaren Wissensquellen verknüpft.

\subsection{Funktionsweise von LLMs}
Die Funktionsweise von LLMs beruht auf wahrscheinlichkeitsbasierten Vorhersagen des jeweils nächsten Tokens auf Basis des bisherigen Kontexts. Dabei werden Wahrscheinlichkeitsverteilungen über mögliche Fortsetzungen berechnet, aus denen sequenziell Text generiert wird. Dieses Verfahren führt zu formal kohärenten und grammatikalisch korrekten Ausgaben, impliziert jedoch kein semantisches Verständnis im menschlichen Sinne.
LLMs reproduzieren das was sie aus den Trainingsdaten gelernt haben und sind daher anfällig für die Übernahme komplexer Strukturen oder Annahmen, die der Zielsetzung verständlicher Texte entgegen wirken können.

\subsection{Prompting-Grundlagen}
Prompting bezeichnet die gezielte Steuerung von Sprachmodellen durch Eingabetexte, die Aufgabenstellung, Stil und gewünschte Eigenschaften der Ausgabe beschreiben.
Für die Textvereinfachung und Leichte Sprache können Prompts genutzt werden, um Anforderungen wie kurze Sätze, einfache Wörter oder explizite Erklärungen vorzugeben. Studien zeigen jedoch, dass Prompting allein keine verlässliche Einhaltung formaler Verständlichkeitsregeln garantiert, da die Umsetzung stark von Modellarchitektur, Trainingsdaten und Kontext abhängt.
Man kann die gewollte Struktur der Antworten zusätzlich mit One- oder Few-shot Prompts weiter verbessern, indem man dem LLM ein oder mehrere Beispiele einer Antwort im Prompt zu Verfügung stellt.

\subsection{Stärken und Schwächen im Kontext Leichte Sprache}
Zu den Stärken von LLMs im Kontext Leichter Sprache zählen ihre Fähigkeit zur schnellen Textproduktion, zur Umformulierung bestehender Texte und zur Generierung inhaltlich verständlicher Vereinfachungen.
Gleichzeitig bestehen deutliche Schwächen, insbesondere die fehlende explizite Ausrichtung auf die Zielgruppe, die unzuverlässige Einhaltung fester Regeln sowie die Gefahr inhaltlicher Abweichungen oder Informationsverluste.
Für einen verlässlichen Einsatz in der leichten Sprache ist daher eine Einbettung von einem LLM in ein System erforderlich, welches die generierten Texte mittels regelbasierter Prüfungen und quantitativer Verständlichkeitsmetriken bewertet und wenn nötig iterativ anpasst \cite[S. 95-99]{maas_handbuch_2018}.